% =====================================================================
%  FICHE 11 — THEME 2 — Corrigé DIFFICILE — Oxydant, réducteur, demi-équations
% =====================================================================
\documentclass[11pt,a4paper]{article}
\usepackage[utf8]{inputenc}
\usepackage[T1]{fontenc}
\usepackage{lmodern}
\usepackage[french]{babel}
\usepackage{amsmath,amssymb}
\usepackage[version=4]{mhchem}
\usepackage{siunitx}
\sisetup{output-decimal-marker={,},per-mode=symbol}
\usepackage{xcolor}
\usepackage{array}
\usepackage{booktabs}
\usepackage{enumitem}
\usepackage[most]{tcolorbox}
\usepackage[a4paper,margin=2.1cm,headheight=26pt,headsep=10pt]{geometry}
\usepackage{fancyhdr}
\usepackage[hidelinks]{hyperref}

\definecolor{primary}{RGB}{146,95,160}
\definecolor{primarydark}{RGB}{92,52,104}
\definecolor{corrcol}{RGB}{0,135,90}
\newcommand{\NO}{\ensuremath{\mathrm{N.O.}}}

\newcounter{exo}
\newtcolorbox{corr}[1][]{enhanced,breakable,boxrule=0.9pt,arc=2pt,colback=corrcol!4,
  colframe=corrcol,fonttitle=\bfseries,coltitle=white,left=3mm,right=3mm,top=2.2mm,bottom=2.2mm,
  title={Corrigé~\refstepcounter{exo}\theexo\ifx\relax#1\relax\relax\else\ \textmd{--- \itshape #1}\fi}}

\pagestyle{fancy}\fancyhf{}
\fancyhead[L]{\small\textcolor{primarydark}{\textbf{Fiche 11 · Thème 2} — Corrigé}}
\fancyhead[R]{\small\textcolor{primarydark}{\textsc{Difficile}}}
\fancyfoot[C]{\small\thepage}
\renewcommand{\headrulewidth}{0.8pt}
\renewcommand{\headrule}{\hbox to\headwidth{\color{primary}\leaders\hrule height \headrulewidth\hfill}}

\begin{document}
\begin{center}
{\LARGE\bfseries\color{primarydark}Thème 2 — Corrigé Difficile}\\[-2pt]
{\color{corrcol}\rule{\textwidth}{1pt}}
\end{center}

\begin{corr}[assembler deux demi-équations]
\begin{enumerate}[label=\textbf{\alph*)},leftmargin=2.2em,topsep=2pt,itemsep=2pt]
  \item \ce{A} ($0\to+\mathrm{II}$) perd 2 électrons : $\ce{A -> A^{2+} + 2 e^-}$.
  \item \ce{B} ($0\to-\mathrm{III}$) capte 3 électrons : $\ce{B + 3 e^- -> B^{3-}}$.
  \item La première échange 2 e$^-$, la seconde 3 e$^-$. On égalise au \textbf{PPCM} $=6$ : on multiplie
        la (a) par 3 et la (b) par 2, soit \textbf{6 électrons} de chaque côté.
  \item Addition membre à membre (les 6 e$^-$ s'éliminent) :
        \[
        \ce{3 A + 2 B -> 3 A^{2+} + 2 B^{3-}}.
        \]
        \ce{A} est oxydé : c'est le \textbf{réducteur}. \ce{B} est réduit : c'est l'\textbf{oxydant}.
\end{enumerate}
\emph{(C'est exactement le schéma « A oxydé par B » de la synthèse du livre, p.\,261.)}
\end{corr}

\begin{corr}[le dosage du permanganate par l'acide oxalique]
\begin{enumerate}[label=\textbf{\alph*)},leftmargin=2.2em,topsep=2pt,itemsep=2pt]
  \item \ce{MnO4^-} : $\NO(\ce{Mn})=+\mathrm{VII}$ ; \ce{Mn^{2+}} : $+\mathrm{II}$.
        Un ion \ce{MnO4^-} capte $7-2=\mathbf{5}$ électrons (réduction).
  \item \ce{(COOH)2} neutre : $2(+1)+2x+4(-2)=0 \Rightarrow 2x=6 \Rightarrow \NO(\ce{C})=+\mathrm{III}$.
        Dans \ce{CO2} : $\NO(\ce{C})=+\mathrm{IV}$. Chaque C perd 1 e$^-$, et il y a 2 C par molécule :
        une \ce{(COOH)2} cède $\mathbf{2}$ électrons (oxydation).
  \item Le permanganate \ce{MnO4^-} est réduit : c'est l'\textbf{oxydant}.
        L'acide oxalique \ce{(COOH)2} est oxydé : c'est le \textbf{réducteur}.
  \item Côté Mn : $2\ \ce{MnO4^-}\times 5 = 10$ électrons captés.
        Côté oxalate : $5\ \ce{(COOH)2}\times 2 = 10$ électrons cédés.
        \textbf{Les deux totaux sont égaux (10 = 10)} : les électrons produits par le réducteur sont
        exactement consommés par l'oxydant — c'est la condition d'une équation redox équilibrée.
\end{enumerate}
\end{corr}

\begin{corr}[laquelle n'est PAS une réaction redox ?]
\begin{enumerate}[label=\textbf{\Alph*.},leftmargin=2.2em,topsep=2pt,itemsep=2pt]
  \item K : $0\to+\mathrm{I}$ (oxydé) ; H de l'eau : $+\mathrm{I}\to 0$ dans \ce{H2} (réduit) $\Rightarrow$ \textbf{redox}.
  \item C : $-\mathrm{IV}$ (\ce{CH4}) $\to +\mathrm{II}$ (\ce{CO}) (oxydé) ; O : $0\to-\mathrm{II}$ (réduit) $\Rightarrow$ \textbf{redox}.
  \item O de \ce{H2O2} ($-\mathrm{I}$) se répartit en \ce{H2O} ($-\mathrm{II}$, réduit) et \ce{O2} ($0$, oxydé) :
        c'est une \textbf{dismutation} $\Rightarrow$ \textbf{redox}.
  \item \ce{CO2 + H2O -> H2CO3} : C reste $+\mathrm{IV}$, O reste $-\mathrm{II}$, H reste $+\mathrm{I}$.
        \textbf{Aucun \NO\ ne change} $\Rightarrow$ ce n'est \textbf{pas un redox} (simple hydratation acide-base).
\end{enumerate}
\textbf{Réponse : D} n'est pas une réaction d'oxydo-réduction (annales 2018-Q25).
\end{corr}

\end{document}
